\section{Introducción}
En esta sesión se retoma el Reward Model mediante el modelo de Bradley-Terry, para explicar la relación entre el modelado probabilístico y estadístico, el aprendizaje profundo y la estimación de máxima verosimilitud.

\begin{importantbox}{Relación entre los tres}
\begin{enumerate}
  \item \textbf{Modelado probabilístico y estadístico}: define el proceso de generación de los datos (por ejemplo, el modelo de preferencias de Bradley-Terry)
  \item \textbf{Estimación de máxima verosimilitud (MLE)}: define la función de pérdida a partir del modelo estadístico
  \item \textbf{Aprendizaje profundo}: sirve como herramienta de resolución, optimizando la pérdida mediante descenso de gradiente
\end{enumerate}
\end{importantbox}

